%%%%%%%%%%  scan idx 124  |  folio 111  %%%%%%%%%%
\SGAPrintPage{111}%
\uline{the derived functor of the tensor product} \((E',E)\longmapsto E'\otimes_{f^{-1}(A)} f^{-1}(E)\)\uline{.} \uline{Let} \(X\in\Ob \underline{S}\)\uline{,} \(X'=f^{-1}(X)\)\uline{,} \(E'\in\Ob\Dpx^-(_{A'}X'_A)\)\uline{,} \(E\in\Ob\Dpx^-(_A X)\)\uline{. Then, when} \(E'\) \uline{is} \(m'\)\uline{-pseudo-coherent (as an object of} \(\Dpx^-(_{A'}\underline{S}')\)\uline{) and acyclic in degree} \(>n'\) \uline{and} \(E\) \uline{is} \(m\)\uline{-pseudo-coherent and acyclic in degree} \(>n\)\uline{,} \(E'\otimes_A^{\mathrm L}E\) \uline{is} \(\sup(m+n',m'+n)\)\uline{-pseudo-coherent and acyclic in degree} \(>n+n'\)\uline{.}

\uline{Proof.} Replacing \(\underline{S}\) by \(\underline{S}_{/X}\), we can suppose that \(X=S\), and we apply (2.13) with \(u=E'\otimes_A^{\mathrm L}:\Dpx^-(_A \underline{S})\to\Dpx^-(_{A'}\underline{S}')\). Let \(U\in\Ob \underline{S}\), \(U'=f^{-1}(U)\), and \(E\) be a complex of left \(A_U\)-Modules. When \(E\) is acyclic in degree \(>n\), \(E'\otimes_A^{\mathrm L}E\) is acyclic in degree \(>n+n'\) (definition of \(\otimes_A^{\mathrm L}\)); on the other hand, for \(E=(A_U)^p\), we have \(E'\otimes_A^{\mathrm L}E=(E'|U')^p\), which is \(m'\)-pseudo-coherent by hypothesis (and (2.12)). The conditions of (2.13) are therefore satisfied, with \(a=m'\), \(b=n'\), whence the conclusion.

In particular:

\sgaStableTarget{sga6.e01.cor.2-16-1}\textbf{Corollary 2.16.1.} \uline{Under the hypotheses of} (2.16), \uline{we have}:

a) \uline{The functor} \(\mathbb{L}f^*: \Dpx^-(_A \underline{S})\to \Dpx^-(_{A'}\underline{S}')\) \uline{(which is simply the functor} \(A'\otimes_A^{\mathrm L}\)\uline{) induces a functor}
\[
  \mathbb{L}f^*: \Dpx^-(_A \underline{S})_{\coh}\longrightarrow \Dpx^-(_{A'}\underline{S}')_{\coh}
\]
(\uline{notation of} (2.6)).

b) \uline{When} \(A\) \uline{is commutative, the functor} \(\otimes_A^{\mathrm L}:\Dpx^-(\underline{S})\times_{\underline{S}}\Dpx^-(\underline{S})\longrightarrow \Dpx^-(\underline{S})\) \uline{induces a functor}
\[
  \otimes_A^{\mathrm L}:\Dpx^-(\underline{S})_{\coh}\times_{\underline{S}}\Dpx^-(\underline{S})_{\coh}
  \longrightarrow \Dpx^-(\underline{S})_{\coh} .
\]

%%%%%%%%%%  scan idx 125  |  folio 112  %%%%%%%%%%
\SGAPrintPage{112}%
\sgaStableTarget{sga6.e01.exercise.2-16-2}\textbf{Exercise 2.16.2.} Let \(\underline{S}\) be a topos, and \(A\to A'\) a morphism of Rings of \(\underline{S}\). We make one of the following hypotheses:
\begin{enumerate}[label=(\roman*)]
\item \(A\to A'\) is surjective with kernel a nilpotent Ideal
\item the functor \(A'\otimes_A-\) on the category of left \(A\)-modules is exact and faithful.
\end{enumerate}
Show that \(E\in\Ob\Dpx^-(_A \underline{S})\) is \(n\)-pseudo-coherent (resp.\ pseudo-coherent) if and only if \(A'\otimes_A^{\mathrm L} E\in\Ob\Dpx^-(_{A'}\underline{S})\) is \(n\)-pseudo-coherent (resp.\ pseudo-coherent).

Hints: First show that a left \(A\)-Module \(L\) is of finite type if and only if \(A'\otimes_A L\) is (cf.\ Bourbaki, Alg. commutative, chap.\ I, \S\ 3, prop.\ 11 and chap.\ II, \S\ 3, prop.\ 4), then reduce to this case by truncation, using (2.5) and (2.10 b)).

\sgaStableTarget{sga6.e01.sec.3}\section*{3. Link with the classical notion of coherence}

\sgaStableTarget{sga6.e01.para.3-0}3.0. The hypotheses and notation are those of (2.0), except that \(\underline{C}_0\) is no longer assumed locally quasi-stable under kernels of epimorphisms.

\sgaStableTarget{sga6.e01.def.3-1}\textbf{Definition 3.1.} \uline{Let} \(F\in\Ob\underline{C}_S\)\uline{. We shall say that} \(F\) \uline{is} precoherent \uline{(relative to} \(\underline{C}_0\)\uline{)}\SGAAuthorFootnote{*}{We introduce this notion here as a technical intermediate step; the only notion to be retained is that of coherence.} \uline{when, for every} \(U\) \uline{over} \(S\) \uline{and every morphism} \(u:E\to F|U\) \uline{with} \(E\in\Ob\underline{C}_{0U}\)\uline{,} \(\Ker(u)\) \uline{is of} \(\underline{C}_0\)\uline{-finite type} (1.2)\uline{. We shall say that} \(F\) \uline{is} coherent \uline{(relative to} \(\underline{C}_0\)\uline{) when} \(F\) \uline{is precoherent and of finite type.}

\sgaStableTarget{sga6.e01.rem.3-2}\textbf{Remark 3.2.} Let \(F\in\Ob\underline{C}_S\). It is clear from the definitions that, when \(F\) is precoherent (resp.\ of finite type), every subobject (resp.\ quotient) of \(F\) is precoherent (resp.\ of finite type).

\sgaStableTarget{sga6.e01.lem.3-3}\textbf{Lemma 3.3.} a) \uline{Let}
\[
  0\longrightarrow F\longrightarrow G\longrightarrow H
\]
%%%%%%%%%%  scan idx 126  |  folio 113  %%%%%%%%%%
\SGAPrintPage{113}%
\uline{be an exact sequence of} \(\underline{C}_S\)\uline{. When} \(G\) \uline{is of finite type and} \(H\) \uline{precoherent,} \(F\) \uline{is of finite type.}

b) \uline{Let}
\[
  F\longrightarrow G\longrightarrow H\longrightarrow 0
\]
\uline{be an exact sequence of} \(\underline{C}_S\)\uline{. When} \(F\) \uline{is of finite type and} \(G\) \uline{precoherent,} \(H\) \uline{is precoherent.}

c) \uline{Let}
\[
  0\longrightarrow F\longrightarrow G\longrightarrow H\longrightarrow 0
\]
\uline{be an exact sequence of} \(\underline{C}_S\)\uline{. When} \(F\) \uline{and} \(H\) \uline{are precoherent (resp.\ of finite type),} \(G\) \uline{is too.}

\uline{Proof.} a) After localizing, we have an epimorphism \(G'\to G\) with \(G'\in\Ob\underline{C}_0\). Hence a commutative diagram with exact rows
\[
\begin{tikzcd}
0\arrow[r]&F'\arrow[r]\arrow[d]&G'\arrow[r]\arrow[d]&H\arrow[d,"\Id"]\\
0\arrow[r]&F\arrow[r]&G\arrow[r]&H .
\end{tikzcd}
\]
\(F'\) is of finite type because \(H\) is precoherent, but \(F'\to F\) is an epimorphism, so \(F\) is of finite type.

b) Let \(U\) be over \(S\), and \(H'\xrightarrow{f}H|U\) a morphism of \(\underline{C}_U\) with \(H'\in\Ob\underline{C}_{0U}\). Changing notation, we can suppose that \(U=S\). By the quasi-liftability of \(\underline{C}_0\) in \(\underline{C}\), there exists (after localizing) a commutative square
\[
\begin{tikzcd}
G' \arrow[r,"p"] \arrow[d] & H' \arrow[d,"f"]\\
G \arrow[r] & H,
\end{tikzcd}
\]
%%%%%%%%%%  scan idx 127  |  folio 114  %%%%%%%%%%
\SGAPrintPage{114}%
where the horizontal arrows are epimorphisms and \(G'\in\Ob\underline{C}_0\). Since \(\Ker(f)\) is a quotient of \(\Ker(fp)\), it suffices to prove that \(\Ker(fp)\) is of finite type, so we can suppose that \(H'=G'\). The arrows \(F\to G\) and \(H'\to G\) define \(g:F\oplus H'\to G\). Putting \(K=\Ker(G\to H)\), we therefore have a commutative diagram in which the rows are exact and \(q\) is an epimorphism:
\[
\begin{tikzcd}
0\arrow[r]&F\arrow[r]\arrow[d,"q"']&F\oplus H'\arrow[r]\arrow[d,"g"]&H'\arrow[r]\arrow[d,"f"]&0\\
0\arrow[r]&K\arrow[r]&G\arrow[r]&H\arrow[r]&0 .
\end{tikzcd}
\]
The snake diagram shows that \(\Ker(g)\to\Ker(f)\) is an epimorphism. But, since \(F\oplus H'\) is of finite type and \(G\) precoherent, \(\Ker(g)\) is of finite type by a), and we are done.

c) Let us first prove the assertion without the ``resp.'' alternative. Up to a change of notation, we must show that, for a morphism \(g:G'\to G\) of \(\underline{C}_S\) with \(G'\in\Ob\underline{C}_0\), \(\Ker(g)\) is of finite type. Denoting by \(H'\) the image of the composite arrow \(G'\to H\), we have a commutative diagram whose rows are exact
\[
\begin{tikzcd}
0\arrow[r]&F'\arrow[r]\arrow[d,"f"']&G'\arrow[r]\arrow[d,"g"]&H'\arrow[r]\arrow[d,"h"]&0\\
0\arrow[r]&F\arrow[r]&G\arrow[r]&H\arrow[r]&0 .
\end{tikzcd}
\]
Since \(h\) is a monomorphism, the snake diagram gives an isomorphism \(\Ker(f)\qis\Ker(g)\). But \(F'\) is of finite type by a), so (again by a)) \(\Ker(f)\) is of finite type, which proves the assertion without the ``resp.'' alternative. Let us pass to the ``resp.'' assertion (strictly speaking, we cannot invoke (2.5 b)) because we do not assume that \(\underline{C}_0\) is locally quasi-stable under kernels of epimorphisms). After localizing, there exist epimorphisms \(F'\xrightarrow{u}F\), \(H'\xrightarrow{v}H\), with
%%%%%%%%%%  scan idx 128  |  folio 115  %%%%%%%%%%
\SGAPrintPage{115}%
\(F'\in\Ob\underline{C}_0\), \(H'\in\Ob\underline{C}_0\). After localizing further, we can suppose that \(v\) factors through \(G\), whence a commutative diagram with exact rows
\[
\begin{tikzcd}
0\arrow[r]&F'\arrow[r]\arrow[d,"u"']&F'\oplus H'\arrow[r]\arrow[d,"w"]&H'\arrow[r]\arrow[d,dashed,"v"]&0\\
0\arrow[r]&F\arrow[r]&G\arrow[r]&H\arrow[r]&0 .
\end{tikzcd}
\]
Since \(u\) and \(v\) are epimorphisms, so is \(w\), and we are done. This completes the proof of (3.3).

We deduce immediately the

\sgaStableTarget{sga6.e01.cor.3-4}\textbf{Corollary 3.4.} \uline{The (strictly) full subcategory of} \(\underline{C}_S\) \uline{formed by the coherent objects is stable under finite projective (resp.\ inductive) limits. In particular, the kernel, cokernel and image of an arrow between two coherent objects are coherent. Let}
\[
  E^0\longrightarrow E^1\longrightarrow E^2\longrightarrow E^3\longrightarrow E^4
\]
\uline{be an exact sequence of} \(\underline{C}_S\)\uline{. When} \(E^0,E^1,E^3,E^4\) \uline{are coherent,} \(E^2\) \uline{is too.}

\sgaStableTarget{sga6.e01.cor.3-5}\textbf{Corollary 3.5.} \uline{Assume that the objects of} \(\underline{C}_0\) \uline{are} \(\underline{C}_0\)\uline{-coherent (which implies in particular that} \(\underline{C}_0\) \uline{is locally quasi-stable under kernels of epimorphisms,} i.e.\ \uline{that the conditions of} (2.0) \uline{hold}).

a) \uline{Equivalent conditions on} \(F\in\Ob\underline{C}_S\):
\begin{enumerate}[label=(\roman*)]
\item \(F\) \uline{is of finite presentation};
\item \(F\) \uline{is coherent};
\item \(F\) \uline{is of} \(\infty\)\uline{-finite presentation} (2.8).
\end{enumerate}

b) \uline{Let} \(F\in\Ob\Dpx^-(\underline{C}_S)\), \(n\in\mathbb Z\). \uline{We have}:
\[
\begin{split}
F\ n\text{-\uline{pseudo-coherent}}
\quad\Longleftrightarrow\quad&
\Ho^i(F)\text{ \uline{coherent for} }i\geq n+1,\\
&\text{\uline{and} }\Ho^n(F)\text{ \uline{of finite type}.}
\end{split}
\]
